% ==============================================================================
% Capítulo 2: Fundamentação teórica e tecnologias
% ==============================================================================
\chapter{Fundamentação teórica e tecnologias}
\label{cap:fundamentacao}

\begin{instrucao}
\orientacao{O que vai na Fundamentação}
Em um trabalho tecnológico, este capítulo reúne dois tipos de conteúdo:
\begin{itemize}
  \item os \textbf{conceitos} necessários para entender o desenvolvimento
        (processos, técnicas de requisitos, padrões de arquitetura, critérios
        de qualidade e de usabilidade), com citações de livros e artigos;
  \item as \textbf{tecnologias} usadas (linguagens, arcabouços, bancos de
        dados, serviços), com uma descrição curta e a justificativa da escolha.
\end{itemize}
Não transforme o capítulo em tutorial. A banca não precisa saber instalar o
Django, mas precisa entender por que ele foi escolhido e que papel cumpre no
sistema. Inclua apenas o que será usado nos capítulos seguintes. Como
referência, de oito a doze páginas costumam ser suficientes.

Escreva com suas palavras e cite as fontes. Copiar texto sem citar é plágio e,
pelas regras do TCC, leva à reprovação com nota zero. Documentação oficial de
tecnologias pode ser citada como página web (veja a entrada
\verb|django2026documentation| em \verb|referencias.bib|).
\end{instrucao}

Este capítulo apresenta os conceitos e as tecnologias que fundamentam o
desenvolvimento do sistema. A \cref{sec:fundamentacao:requisitos} trata da
especificação de requisitos com histórias de usuário. A
\cref{sec:fundamentacao:arquitetura} apresenta o padrão de arquitetura
adotado, a \cref{sec:fundamentacao:usabilidade} discute a avaliação de
usabilidade e a \cref{sec:fundamentacao:tecnologias} descreve as tecnologias
utilizadas.

% ------------------------------------------------------------------------------
\section{Requisitos e histórias de usuário}
\label{sec:fundamentacao:requisitos}

Requisitos descrevem o que um sistema deve fazer e as restrições sob as quais
deve operar \cite{sommerville2018engenharia}. Eles costumam ser divididos em
\acp{rf}, que descrevem serviços e comportamentos do sistema, e \acp{rnf}, que
tratam de propriedades como desempenho, segurança, usabilidade e
manutenibilidade \cite{sommerville2018engenharia}.

Em processos ágeis, os requisitos funcionais são frequentemente registrados
como histórias de usuário, frases curtas escritas do ponto de vista de quem
usará o sistema, no formato ``como \emph{papel}, quero \emph{ação} para
\emph{benefício}'' \cite{cohn2004user}. Cada história é acompanhada de
critérios de aceitação, que definem quando ela está concluída.
\textcite{cohn2004user} recomenda que boas histórias sejam independentes,
negociáveis, valiosas para o usuário, estimáveis, pequenas e testáveis.

\begin{instrucao}
\orientacao{Citação direta}
A citação direta reproduz literalmente o texto da fonte e sempre informa a
página. Com até três linhas, fica no próprio parágrafo, entre aspas duplas.
\begin{exemplo}
A portaria define alucinação como
``informação aparentemente plausível,
porém incorreta, imprecisa ou
inventada'' \cite[p.~1]{ufes2026portaria}.
\end{exemplo}
Com mais de três linhas, use o ambiente \verb|citacao|, que aplica recuo de
4~cm, fonte menor e espaçamento simples, sem aspas.
\begin{exemplo}
\begin{citacao}
O usuário é responsável pela precisão,
pela originalidade e pela adequação dos
conteúdos gerados ou modificados por
sistema de inteligência artificial que
utilizar, devendo verificá-los antes de
sua utilização, considerada a
possibilidade de alucinação
\cite[p.~5]{ufes2026portaria}.
\end{citacao}
\end{exemplo}
Use citações diretas com moderação. Em uma Fundamentação bem escrita, a maior
parte das citações é indireta (paráfrase com indicação da fonte).
\end{instrucao}

% ------------------------------------------------------------------------------
\section{Arquitetura de software e o padrão MTV}
\label{sec:fundamentacao:arquitetura}

A arquitetura de um sistema é o conjunto de estruturas necessárias para
raciocinar sobre ele, formado pelos elementos de software, pelas relações
entre eles e pelas propriedades de ambos \cite{bass2021software}. Em
aplicações web, um padrão recorrente é separar a lógica de dados, a lógica de
apresentação e o tratamento das requisições.

O arcabouço Django organiza as aplicações segundo o padrão \ac{mtv}
\cite{django2026documentation}. O modelo define a estrutura dos dados e as
regras de negócio associadas a eles e é mapeado para tabelas do banco de dados
por meio de um \ac{orm}. A visão recebe a requisição, consulta os modelos e
escolhe a resposta. O \emph{template} define como os dados são apresentados
ao usuário. A \cref{fig:mtv} mostra o caminho de uma requisição nesse padrão.

\begin{figure}[htbp]
  \centering
  \caption{Fluxo de uma requisição no padrão \acs{mtv} do Django}
  \label{fig:mtv}
  \begin{tikzpicture}[
      node distance=8mm and 14mm,
      caixa/.style={draw, rounded corners=2pt, minimum width=2.6cm,
                    minimum height=9mm, align=center, font=\footnotesize},
      seta/.style={-{Stealth}, thick},
    ]
    \node[caixa] (nav) {Navegador};
    \node[caixa, right=of nav] (url) {Roteador de URLs\\(\texttt{urls.py})};
    \node[caixa, right=of url] (visao) {Visão\\(\texttt{views.py})};
    \node[caixa, below=of visao] (modelo) {Modelo\\(\texttt{models.py})};
    \node[caixa, below=of modelo, fill=gray!10] (bd) {Banco de dados};
    \node[caixa, below=of url] (tpl) {\emph{Template}\\(HTML)};
    \draw[seta] (nav) -- node[above, font=\scriptsize] {requisição} (url);
    \draw[seta] (url) -- (visao);
    \draw[seta] (visao) -- (modelo);
    \draw[seta] (modelo) -- node[right, font=\scriptsize] {\acs{orm}} (bd);
    \draw[seta] (visao) -- (tpl);
    \draw[seta] (tpl.west) -| node[below, pos=0.25, font=\scriptsize] {resposta} (nav.south);
  \end{tikzpicture}
  \fonte{Elaborado pelo próprio autor com base em \textcite{django2026documentation}.}
\end{figure}

\begin{instrucao}
\orientacao{Figuras: inclusão, citação e qualidade}
Pela NBR 14724, a identificação da figura (a palavra ``Figura'', o número e
o título) fica \textbf{acima} dela e a fonte fica \textbf{abaixo}, mesmo
quando a figura foi feita por você. Por isso, no LaTeX, o \verb|\caption| vem
antes do conteúdo e o \verb|\fonte| depois. Para material próprio, use
\verb|\fonteautor|. Para figura baseada em outra obra, escreva, por exemplo,
\verb|\fonte{Adaptado de \textcite{chave}.}|.
\begin{exemplocodigo}
\begin{figure}[htbp]
  \centering
  \caption{Tela de reserva de horários}
  \label{fig:tela-reserva}
  \includegraphics[width=0.9\textwidth]{figuras/tecnologico/tela-reserva.png}
  \fonteautor
\end{figure}
\end{exemplocodigo}

\textbf{Toda figura, tabela, quadro ou código incluído no trabalho deve ser
citado no texto pelo menos uma vez}, de preferência antes de aparecer. Ao
citar, escreva o nome com inicial maiúscula seguido do número, como em ``a
Figura 3 mostra'', e não ``a figura 3''. O mesmo vale para Tabela, Quadro,
Código, Equação e Seção. Os comandos \verb|\cref| e \verb|\Cref| já escrevem o
nome dessa forma. Não basta citar. Explique no texto o que deve ser observado. Diagramas de
arquitetura, de classes e de sequência pedem explicação no texto, elemento por
elemento ou pelo menos das partes principais.

\textbf{Como obter figuras de boa qualidade}
\begin{itemize}
  \item Prefira figuras \textbf{vetoriais} (PDF, EPS ou desenhadas em TikZ),
        que não perdem nitidez com o zoom nem na impressão. Esta figura foi
        desenhada em TikZ diretamente no texto.
  \item Diagramas UML feitos no draw.io (diagrams.net), no PlantUML, no Astah
        ou em ferramentas semelhantes devem ser exportados em PDF ou SVG
        (convertido para PDF), nunca capturados com \emph{print screen}.
  \item Capturas de tela do sistema devem ser salvas em PNG, na resolução
        original e com o zoom do navegador em 100\%. Recorte apenas a área de
        interesse e evite JPEG, porque a compressão borra textos e linhas.
  \item Ajuste apenas a largura (\verb|width=|). Alterar largura e altura ao
        mesmo tempo distorce a imagem.
  \item O texto dentro da figura deve ter tamanho próximo ao das legendas.
        Se for preciso dar zoom para ler, a figura precisa ser refeita ou
        dividida.
  \item Use cores que continuem distinguíveis em escala de cinza e para
        pessoas com daltonismo.
  \item Mantenha o mesmo estilo e o mesmo alinhamento (centralizado) em todas
        as figuras do trabalho, como pede a NBR 14724:2024.
\end{itemize}
\end{instrucao}

% ------------------------------------------------------------------------------
\section{Usabilidade e a escala SUS}
\label{sec:fundamentacao:usabilidade}

A usabilidade é a medida em que um sistema pode ser usado por usuários
específicos para atingir objetivos com eficácia, eficiência e satisfação em um
contexto de uso \cite{iso2018usability}. Um dos instrumentos mais usados para medir a satisfação é a
\ac{sus}, um questionário de dez afirmações respondidas em escala de cinco
pontos que resulta em uma pontuação de 0 a 100 \cite{brooke1996sus}.
\textcite{bangor2009determining} associaram faixas dessa pontuação a adjetivos,
o que facilita a interpretação. Pontuações em torno de 71 correspondem ao
adjetivo ``bom''. Pontuações em torno de 86 correspondem a ``excelente''.

% ------------------------------------------------------------------------------
\section{Tecnologias utilizadas}
\label{sec:fundamentacao:tecnologias}

O \cref{qua:tecnologias} resume as tecnologias usadas no desenvolvimento, com
a versão e o papel de cada uma. A escolha do Django se deveu à experiência
prévia do autor com Python, ao sistema de autenticação e de administração já
incluídos no arcabouço e à documentação extensa
\cite{django2026documentation}. O PostgreSQL foi escolhido por ser o banco de
dados já mantido pela equipe de tecnologia do centro, o que facilita a
implantação.

\begin{quadro}[htbp]
  \centering
  \caption{Tecnologias utilizadas no desenvolvimento}
  \label{qua:tecnologias}
  \small
  \begin{tabularx}{\textwidth}{|l|c|X|}
    \hline
    \textbf{Tecnologia} & \textbf{Versão} & \textbf{Papel no sistema} \\
    \hline
    Python & 3.12 & Linguagem de programação do servidor \\
    \hline
    Django & 5.2 & Arcabouço web (rotas, visões, \acs{orm}, autenticação) \\
    \hline
    PostgreSQL & 16 & Banco de dados relacional \\
    \hline
    Bootstrap & 5.3 & Componentes de interface e leiaute responsivo \\
    \hline
    pytest & 8.3 & Execução dos testes automatizados \\
    \hline
    Docker & 27 & Empacotamento e implantação em contêineres \\
    \hline
  \end{tabularx}
  \fonteautor
\end{quadro}
